\section{Limitations and Future Work}
\label{sec:limitations}

A successful conversion does not preserve every aspect of the source.
Some differences follow from the content types and HTML that H5P accepts; others follow from the
converter's treatment of Markdown.
\autoref{tab:limitations} summarizes the main approximations and losses.

\subsection{Approximations and Information Loss}
\label{sec:lm_conversion}

\begin{table}[H]
\centering
\small
\begin{tabular}{@{}>{\raggedright\arraybackslash}p{0.23\columnwidth}>{\raggedright\arraybackslash}p{0.73\columnwidth}@{}}
\toprule
Construct & Limitation or Approximation \\
\midrule
Headings and lists & Headings are shifted down one level and clamped at \code{h3}. Tight lists receive the spacing of loose lists. \\
\addlinespace[2pt]
Quotes and inline formatting & Block quotes and alerts become shaded boxes, inserted text receives a green background, and task list checkboxes become characters. \\
\addlinespace[2pt]
Image placement & Images become separate content blocks, interrupting surrounding lists, quotes, or footnotes. \\
\addlinespace[2pt]
Image size & Image blocks have a width attribute, but it is ignored: images occupy the full width of their column. Markdown image size attributes therefore have no effect. \\
\addlinespace[2pt]
Image formats and availability & Unsupported formats, including SVG, and unavailable images are omitted with a warning. \\
\addlinespace[2pt]
Tables with images & Tables with up to four columns become pseudo-tables. Wider tables retain placeholders and place the images below the table. \\
\addlinespace[2pt]
Footnotes & References use bracketed numbers, and the notes form an ordered list at the end of the chapter. \\
\addlinespace[2pt]
Links & Links to other chapters become textual references. Links to Markdown files outside the book lose their destination. \\
\addlinespace[2pt]
Raw HTML & Tags and styles remain subject to the platform's restrictions. Hand-written anchors and links are not rewritten and may therefore collide with generated ids or target the wrong frame. \\
\addlinespace[2pt]
Code and questions & Unknown code languages remain uncolored. Malformed question blocks fall back to code with a warning. \\
\addlinespace[2pt]
Formulas in tables & An unescaped vertical bar splits a table cell before formula parsing. Authors must write \code{\textbackslash|} or \code{\textbackslash{}mid} instead. \\
\addlinespace[2pt]
Metadata & YAML front matter is discarded, and chapter titles derived from file names are not capitalized automatically. \\
\bottomrule
\end{tabular}
\caption[Approximations and information loss]{The main approximations and information loss during conversion. Warnings identify unsupported input, but visual inspection in the player remains necessary.}
\label{tab:limitations}
\end{table}

\subsection{Deployment and Merge Limitations}
\label{sec:lm_platform}

Formulas require the \code{H5P.MathDisplay} add-on on the target site.
The first upload also requires permission to install any bundled libraries that the site does not
already provide.
These requirements cannot be satisfied by the package alone.

Beyond these deployment requirements, the merge can misplace additions when generated blocks are reordered or closely resemble one another.
In addition, generated blocks deleted in H5P reappear after merging, since Markdown remains their source.
Library versions are adjusted without running the libraries' migrations, so these changes produce a
warning.

\subsection{Future Work}
\label{sec:lm_future}

The question syntax could be extended with further question types,
such as fill-in-the-blank or true/false questions.
Each would need a readable fence syntax and a renderer for the corresponding H5P library.
The existing placement mechanism could then handle the resulting blocks.

Packaging could also include only the libraries a book actually uses.
This would reduce package size and the repeated archive work,
but the merge would also need to account for the changing set of libraries.

Finally, automatic rebuilds could run the CLI in a repository's CI pipeline and publish the resulting
package as a build artifact after each change.
Uploading it to Moodle, however, would remain a manual step.
